\scanpage{59}
% manuscript line 1
\uline{2. Produits extérieurs hilbertiens. Applications : des inégalités diverses}

% manuscript line 2
Soient $E_1,\ldots,E_n$ $n$ espaces de Hilbert. Alors on peut
encore, sur leur produit tensoriel $\bigotimes_{1\leq i\leq n}E_i$, introduire
un produit scalaire, donné sur les tenseurs décomposables
par

% manuscript line 6
\[
(a_1\otimes\cdots\otimes a_n,b_1\otimes\cdots\otimes b_n)=(a_1,b_1)\cdots(a_n,b_n) \tag*{(1)}
\]

% manuscript line 7
On vérifie comme au N° 1 qu'on obtient là une
forme sesquilinéaire \deleted{p\unsure{os}} strictement positive, \deleted{et que} d'où un
\deleted{$\unsure{\sum_{\alpha}\lambda_\alpha E_\alpha|\bigotimes_{i\in I_\alpha}^{(2)}}$ est une \unsure{base orthonormale, dont}}
\deleted{le complété de Hilbert \unsure{s'obtient} \illegible{1} noté}
\inserted{\uline{appelé \textquotedblleft produit tensoriel hilbertien des $E_i$\textquotedblright}, noté}
$\overset{(2)}{\bigotimes}_{1\leq i\leq n}E_i$, ou $\overset{n}{\otimes}\vphantom{\otimes}^{(2)}E$ si tous les $E_i$ se réduisent
à un $=$ espace $E$. Si, pour \deleted{\unsure{tout}} $i$,
$(e^i_{\alpha_i})_{\alpha_i\in A_i}$ est une base orthonormale dans $E_i$,
alors on obtient une base orthonormale
$(e_\alpha)_{\alpha\in A}$ de $\overset{(2)}{\bigotimes}_{1\leq i\leq n}E_i$, où $A=A_1\times\cdots\times A_n$,
et pour $\alpha=(\alpha_1,\ldots,\alpha_n)$, on pose $e_\alpha=e^1_{\alpha_1}\otimes\cdots e^n_{\alpha_n}$.

% manuscript line 9a
% margin note written sideways
\margin{on a un isomorphisme canonique, noté encore $\sigma$, de $\overset{(2)}{\bigotimes}_{1\leq i\leq n}E_i$ sur $\overset{(2)}{\bigotimes}_{1\leq i\leq n}E_{\sigma\cdot i}$, \deleted{\unsure{isomorphisme} \illegible{2}}. Sur la formule (1), on voit aussitôt \uline{l'associativité} de l'opération du produit tensoriel hilbertien : si $1\leq\alpha<\beta\ldots<\gamma<n$, on a un isomorphisme canonique de $\overset{\unsure{(2)|n}}{\bigotimes}_{1\leq i\leq n}E_i$ avec $(\overset{(2)}{\bigotimes}_{1\leq i\leq\alpha}E_i)\overset{(2)}{\otimes}(\overset{(2)}{\bigotimes}_{\alpha<i\leq\beta}E_i)\overset{(2)}{\otimes}\cdots\overset{(2)}{\otimes}(\overset{(2)}{\bigotimes}_{\gamma<i\leq n}E_i)$. On voit de même qu'une \deleted{\unsure{autre}} permutation $\sigma\in\mathfrak S_n$, sur les tenseurs décomposables par $\sigma(x_1\otimes\cdots\otimes x_n)=x_{\sigma1}\otimes\cdots\otimes x_{\sigma n}$ (commutativité du produit tensoriel hilbertien).}

% manuscript line 18
\deleted{Soit une} \inserted{Supposons donnée} une \deleted{\unsure{2ième}} deuxième \deleted{séquence}
$(F_i)_{1\leq i\leq n}$ d'espaces de Hilbert, et soit, pour \deleted{\unsure{tout}} $i$, \deleted{$A$} $u_i$ une application
lin. continue de $E_i$ dans $F_i$, alors l'application
$\bigotimes_{1\leq i\leq n}u_i$ de $\bigotimes E_i$ dans $\bigotimes F_i$ est continue
pour leurs normes hilbertiennes, et \uline{de façon précise}

% manuscript line 23
\[
\|u_1\otimes\cdots\otimes u_n\|\leq\|u_1\|\,\deleted{\unsure{\otimes}}\cdots\|u_n\| \tag*{(2)}
\]

% manuscript line 24
Il suffit en effet de le démontrer quand tous les $u_i$
sauf l'un sont l'identité. \deleted{On trouve le résultat en écrivant}

% manuscript line 26
% archive sheet number: 60
